\section{주분석 결과}
\label{sec:main-results}

이 절은 사전에 고정한 두 공동 주패널을 같은 위계로 보고한다. Panel A는
삼쌍승 배당상한이 없는 공통 clean 표본의 점가격 비교이고, Panel B는 전체
분석표본에서 반올림과 상한이 허용하는 가격집합을 그대로 유지한
부분식별 결과다. 모든 수치는 재현 스크립트가 생성한 표를 직접 불러온다.

\subsection{Clean 표본의 선택과 구성}

\begin{table}[htbp]
\centering
\caption{Clean 표본과 상한 포함 표본의 구성 비교}
\label{tab:main-selection}
\resizebox{\textwidth}{!}{\begin{tabular}{lrrrrrrrr}
\toprule
표본 & 경주 수 & 출전두수 중앙값 [IQR] & 2016--19 비중 & 서울(1) & 제주(2) & 부산경남(3) & 비검열 배당 95\% & 99\% \\
\midrule
clean & 3,338 & 9 [8, 10] & 36.3\% & 31.4\% & 50.9\% & 17.7\% & 4,625.6 & 6,713.9 \\
상한 포함 & 15,963 & 11 [10, 12] & 52.7\% & 43.6\% & 25.1\% & 31.3\% & 8,464.3 & 9,662.2 \\
\bottomrule
\end{tabular}
}
\end{table}

표~\ref{tab:main-selection}은 clean 표본이 전체 시장의 무작위 축소판이 아님을
보여준다. clean 경주의 출전두수 중앙값은 9두[IQR 8--10]인 반면 상한 조합이
있는 경주는 11두[10--12]다. 2016--2019년의 비중도 각각 36.0\%와 52.7\%로
다르고, 세 경마장 코드의 구성비 역시 눈에 띄게 다르다. 또한 상한 경주의
비검열 삼쌍승 배당 꼬리는 훨씬 두꺼워 95백분위가 8,464.3, 99백분위가
9,662.2인 반면 clean 표본은 각각 4,590.3과 6,629.8이다. 따라서 상한 발생은
출전두수와 가격꼬리의 복잡성과 체계적으로 연관되어 있으며, Panel A만으로 전체
시장을 대표한다고 볼 수 없다.

이 차이가 Panel A와 Panel B의 TV 차이를 인과적으로 설명한다고 주장하지는
않는다. 다만 표본선택 가능성이 실질적이라는 점은 분명하다. clean 표본 안에서
주모형 중앙 TV의 연도별 범위는 쌍승 0.0521--0.0604, 복승
0.0608--0.0698, 삼복승 0.0383--0.0500으로 비교적 안정적이지만, 출전두수와
경마장별로도 일정한 이질성이 남는다. 이 때문에 이하에서는 clean 점추정과
전체표본 부분식별을 반드시 함께 판정한다.

\begin{table}[htbp]
\centering
\caption{Panel A와 Panel B의 이질성 범위 비교}
\label{tab:main-heterogeneity}
\resizebox{\textwidth}{!}{\begin{tabular}{llrrr}
\toprule
승식 & 구분 & Panel A 중앙 TV 범위 & Panel B TV 하한 중앙값 범위 & Panel B TV 상한 중앙값 범위 \\
\midrule
쌍승 & 연도 & [0.0521, 0.0604] & [0.0528, 0.0606] & [0.0625, 0.0789] \\
쌍승 & 경마장 & [0.0504, 0.0608] & [0.0534, 0.0625] & [0.0684, 0.0722] \\
쌍승 & 출전두수 & [0.0500, 0.0666] & [0.0475, 0.0593] & [0.0557, 0.1564] \\
복승 & 연도 & [0.0608, 0.0698] & [0.0587, 0.0674] & [0.0711, 0.0892] \\
복승 & 경마장 & [0.0559, 0.0714] & [0.0599, 0.0715] & [0.0779, 0.0830] \\
복승 & 출전두수 & [0.0536, 0.0766] & [0.0435, 0.0674] & [0.0637, 0.1750] \\
삼복승 & 연도 & [0.0383, 0.0500] & [0.0408, 0.0494] & [0.0496, 0.0666] \\
삼복승 & 경마장 & [0.0389, 0.0493] & [0.0420, 0.0519] & [0.0564, 0.0595] \\
삼복승 & 출전두수 & [0.0363, 0.0471] & [0.0276, 0.0475] & [0.0472, 0.1634] \\
\bottomrule
\end{tabular}
}
\end{table}

표~\ref{tab:main-heterogeneity}은 연도·경마장·출전두수별 이질성을 두 패널에서
같이 비교한다. 연도와 경마장 차원에서는 Panel A의 중앙 TV 범위와 Panel B의
TV 하한 중앙값 범위가 대체로 같은 규모와 순서를 유지한다. 즉 clean 표본에서
보인 상대적 정합성 패턴이 특정 연도나 특정 경마장에만 의존한다는 증거는 약하다.
반면 출전두수별 Panel B의 보수적 TV 상한 범위는 쌍승 0.0557--0.1564,
복승 0.0637--0.1750, 삼복승 0.0472--0.1634로 크게 넓어진다. 이는 일부
출전두수 집단에서 상한·반올림이 허용하는 가격집합 자체가 넓어져 최악경우
불확실성이 커진다는 뜻이지, 해당 집단의 실제 오차가 그 상한값만큼 크다는
뜻은 아니다. 따라서 이질성 차원에서도 Panel B는 Panel A의 방향성을 대체로
지지하지만, 출전두수에 따른 부분식별 폭의 차이를 무시할 수는 없다.

\subsection{Clean 표본의 점가격 비교}

\begin{table}[htbp]
\centering
\caption{Panel A: 삼쌍승 주변가격과 실제 승식 가격의 거리}
\label{tab:main-panel-a}
\resizebox{\textwidth}{!}{\begin{tabular}{llrrrr}
\toprule
승식 & 모형 & 경주 수 & 중앙 TV & 95\% CI 하한 & 95\% CI 상한 \\
\midrule
exacta & harville & 3,321 & 0.1120 & 0.1110 & 0.1130 \\
exacta & main & 3,321 & 0.0555 & 0.0550 & 0.0560 \\
exacta & other\_race & 3,320 & 0.5846 & 0.5780 & 0.5897 \\
exacta & permutation & 3,321 & 0.4807 & 0.4769 & 0.4848 \\
exacta & uniform & 3,321 & 0.4358 & 0.4326 & 0.4385 \\
quinella & harville & 3,321 & 0.1091 & 0.1082 & 0.1104 \\
quinella & main & 3,321 & 0.0636 & 0.0629 & 0.0646 \\
quinella & other\_race & 3,320 & 0.5617 & 0.5564 & 0.5673 \\
quinella & permutation & 3,321 & 0.4353 & 0.4319 & 0.4398 \\
quinella & uniform & 3,321 & 0.4128 & 0.4090 & 0.4155 \\
trio & harville & 3,321 & 0.1519 & 0.1501 & 0.1532 \\
trio & main & 3,321 & 0.0444 & 0.0438 & 0.0449 \\
trio & other\_race & 3,320 & 0.6013 & 0.5970 & 0.6055 \\
trio & permutation & 3,321 & 0.4945 & 0.4917 & 0.4975 \\
trio & uniform & 3,321 & 0.4476 & 0.4445 & 0.4508 \\
win & harville & 3,321 & 0.0000 & 0.0000 & 0.0000 \\
win & main & 3,321 & 0.0679 & 0.0665 & 0.0691 \\
win & other\_race & 3,320 & 0.4649 & 0.4591 & 0.4694 \\
win & permutation & 3,321 & 0.3368 & 0.3337 & 0.3403 \\
win & uniform & 3,321 & 0.3293 & 0.3265 & 0.3316 \\
\bottomrule
\end{tabular}
}
\end{table}

표~\ref{tab:main-panel-a}에서 삼쌍승 상태가격을 실제 경주의 단승·쌍승·복승·
삼복승 결과공간으로 주변화한 주모형은, 단승을 제외한 세 복합승식에서
Harville보다 작은 TV를 보인다. 동일 경주에서 삼쌍승 가격을 무작위로 섞은
순열 기준, 같은 출전두수의 다른 경주를 가져온 기준, 균등분포 기준과의 격차는
더 크다. 이는 높은 적합도가 단순히 상태공간의 차원이나 출전두수만으로 생긴
기계적 결과가 아님을 보여준다.

단승의 Harville 열은 별도의 예측모형 검정으로 해석하지 않는다. Harville의
입력 자체가 같은 경주의 정규화 단승 가격이므로 단승으로 다시 주변화하면
구성상 단승 가격을 재현한다. 단승 열은 삼쌍승 주변가격이 단승 가격과 얼마나
가까운지를 보여주는 교차풀 진단이고, Harville 대비 우위의 의미 있는 검정은
쌍승·복승·삼복승에서 이루어진다.

\subsection{Panel A의 보정회귀와 가중 강건성}

\begin{table}[htbp]
\centering
\caption{Panel A 보정회귀와 가중·군집 강건성}
\label{tab:main-calibration}
\resizebox{\textwidth}{!}{\begin{tabular}{llrrrrr}
\toprule
승식 & 가중 & $\hat\beta$ & 경주군집 SE & 경주--날짜 SE & 경주별 $\beta$ 중앙값 & $[0.9,1.1]$ 비율 \\
\midrule
단승 & 경주균등 & 0.8943 & 0.0016 & 0.0018 & 0.9010 & 48.5\% \\
단승 & 조합균등 & 0.8949 & 0.0016 & 0.0017 & 0.9010 & 48.5\% \\
쌍승 & 경주균등 & 1.0093 & 0.0011 & 0.0013 & 1.0005 & 92.8\% \\
쌍승 & 조합균등 & 1.0084 & 0.0010 & 0.0012 & 1.0005 & 92.8\% \\
복승 & 경주균등 & 0.9560 & 0.0013 & 0.0016 & 0.9536 & 79.7\% \\
복승 & 조합균등 & 0.9611 & 0.0012 & 0.0015 & 0.9536 & 79.7\% \\
삼복승 & 경주균등 & 0.9805 & 0.0009 & 0.0012 & 0.9852 & 94.1\% \\
삼복승 & 조합균등 & 0.9898 & 0.0008 & 0.0010 & 0.9852 & 94.1\% \\
\bottomrule
\end{tabular}
}
\end{table}

식~\eqref{eq:calibration}의 보정기울기를 출전두수 고정효과와 함께 추정한 결과,
경주균등가중 $\hat\beta$는 단승 0.8946, 쌍승 1.0096, 복승 0.9564,
삼복승 0.9807이다. 조합균등가중으로 바꾸어도 각각 0.8952, 1.0087,
0.9616, 0.9901로 큰 변화가 없다. 경주--날짜 이중군집 표준오차는
경주균등가중에서 각각 0.0018, 0.0013, 0.0016, 0.0012로 경주군집
표준오차보다 조금 크지만 결과의 실질적 크기는 바뀌지 않는다. 표본이 크므로
기울기가 1과 통계적으로 정확히 같은지를 주장하기보다, 1에서 떨어진 정도와
경주별 분포를 함께 본다.

경주별 보정기울기의 중앙값은 단승 0.9011, 쌍승 1.0007, 복승 0.9536,
삼복승 0.9852이고, 경주별 기울기가 $[0.9,1.1]$에 드는 비율은 각각
48.6\%, 92.7\%, 79.6\%, 94.1\%다. 따라서 TV가 작은 것과 모든 승식에서
정확한 단위기울기가 성립하는 것은 같은 명제가 아니다. 보정회귀까지 함께 보면
쌍승과 삼복승의 가격정렬이 특히 강하고, 복승은 완만한 기울기 차이가 남으며,
단승은 더 뚜렷한 차이를 보인다. 이 표는 clean Panel A의 설명적 강건성 결과이고
Panel B의 부분식별 판정을 대체하지 않는다.

\subsection{전체표본의 부분식별 결과}

\begin{table}[htbp]
\centering
\caption{Panel B: 전체표본의 TV 부분식별 경계. 주모형·순열·타경주 기준은
게시 배당의 허용 가격집합을 유지하지만, Harville은 게시 단승 점벡터에서
계산한 퇴화한 점예측을 사용하므로 Harville 비교는 대칭적인 부분식별이 아니다.}
\label{tab:main-panel-b}
\resizebox{\textwidth}{!}{\begin{tabular}{llrrrrr}
\toprule
승식 & 모형 & 경주 수 & TV 하한 중앙값 & 하한 95\% CI & TV 상한 중앙값 & 상한 95\% CI \\
\midrule
쌍승 & 단승--Harville & 19,284 & 0.1195 & [0.1190, 0.1199] & 0.1239 & [0.1234, 0.1244] \\
쌍승 & 삼쌍승 주변화 & 19,284 & 0.0567 & [0.0565, 0.0569] & 0.0703 & [0.0700, 0.0706] \\
쌍승 & 타경주 & 19,283 & 0.7012 & [0.6993, 0.7034] & 0.7161 & [0.7140, 0.7187] \\
쌍승 & 동일경주 순열 & 19,284 & 0.5849 & [0.5831, 0.5871] & 0.5985 & [0.5965, 0.6005] \\
쌍승 & 균등 & 19,284 & 0.5395 & [0.5380, 0.5410] & 0.5428 & [0.5411, 0.5446] \\
복승 & 단승--Harville & 19,284 & 0.1155 & [0.1148, 0.1161] & 0.1226 & [0.1219, 0.1232] \\
복승 & 삼쌍승 주변화 & 19,284 & 0.0638 & [0.0635, 0.0641] & 0.0809 & [0.0805, 0.0812] \\
복승 & 타경주 & 19,283 & 0.6766 & [0.6744, 0.6788] & 0.6937 & [0.6916, 0.6961] \\
복승 & 동일경주 순열 & 19,284 & 0.5356 & [0.5338, 0.5373] & 0.5513 & [0.5490, 0.5532] \\
복승 & 균등 & 19,284 & 0.5083 & [0.5066, 0.5101] & 0.5142 & [0.5123, 0.5159] \\
삼복승 & 단승--Harville & 19,284 & 0.1647 & [0.1641, 0.1655] & 0.1687 & [0.1680, 0.1695] \\
삼복승 & 삼쌍승 주변화 & 19,284 & 0.0454 & [0.0452, 0.0455] & 0.0583 & [0.0581, 0.0586] \\
삼복승 & 타경주 & 19,283 & 0.7215 & [0.7196, 0.7233] & 0.7359 & [0.7340, 0.7382] \\
삼복승 & 동일경주 순열 & 19,284 & 0.6153 & [0.6135, 0.6171] & 0.6270 & [0.6251, 0.6292] \\
삼복승 & 균등 & 19,284 & 0.5556 & [0.5537, 0.5572] & 0.5585 & [0.5567, 0.5604] \\
단승 & 단승--Harville & 19,284 & 0.0000 & [0.0000, 0.0000] & 0.0094 & [0.0094, 0.0094] \\
단승 & 삼쌍승 주변화 & 19,284 & 0.0603 & [0.0598, 0.0607] & 0.0862 & [0.0857, 0.0868] \\
단승 & 타경주 & 19,283 & 0.5694 & [0.5673, 0.5715] & 0.5933 & [0.5908, 0.5957] \\
단승 & 동일경주 순열 & 19,284 & 0.4143 & [0.4123, 0.4160] & 0.4368 & [0.4350, 0.4390] \\
단승 & 균등 & 19,284 & 0.4108 & [0.4091, 0.4124] & 0.4241 & [0.4226, 0.4258] \\
\bottomrule
\end{tabular}
}
\end{table}

Panel B에서도 복합승식에 대한 순위는 유지된다. 주모형의 보수적 TV 상한조차
Harville의 TV 하한보다 낮은 경주의 비율은 쌍승 97.6\%, 복승 86.8\%,
삼복승 99.8\%이고, 경주균등가중 중앙 개선폭의 bootstrap 하한도 세 승식 모두
양수다. 순열·타경주·균등 기준에 대해서는 그 차이가 더욱 뚜렷하다. 따라서
Panel A의 결과는 상한 없는 17.2\% 표본에서만 나타난 현상으로 설명하기 어렵다.

그러나 이 결과를 ``완전한 공통 상태가격''의 증거로 읽어서는 안 된다.
사전에 정한 절대 TV 기준 0.05에서는 어느 목표 승식도 Panel A와 Panel B를
동시에 통과하지 않는다. clean 표본에서 가장 작은 절대오차를 보이는 삼복승도
전체표본의 보수적 상한까지 고려하면 이 강한 기준을 인증하지 못한다. 더 엄격한
민감도 기준 0.025에서도 네 목표 승식 모두 두 패널의 공동 판정을 통과하지
못한다. 반면 민감도 기준 0.10에서는 두 패널이 모두 통과한다. 따라서 주결론은
``정확한 등식''보다 ``기준모형을 크게 능가하는 높은 교차풀 정합성, 그러나
무시할 수 없는 잔여 가격차''로 제한한다.

\subsection{기준모형 대비 개선과 순서정보}

\begin{table}[htbp]
\centering
\caption{기준모형 대비 TV 개선폭. 단승--Harville 행은 Harville의 입력 자체가
단승 정규화 가격이어서 구성상 항등식이며, 실질적인 기준모형 우위 검정으로
해석하지 않는다. Panel B의 Harville 행들은 게시 단승 점벡터에서 만든
퇴화한 점예측에 조건부이며 다른 가격집합 기준과 완전히 대칭적인 부분식별 비교가 아니다.
각 95\% 신뢰구간은 개별 구간이며 다중비교 보정을 적용하지 않았다.
표의 $\dagger$는 같은 donor가 여러 대상경주에 재사용되지만 해당 신뢰구간에는
donor 군집보정이 적용되지 않은 보조 기준모형을 뜻한다.}
\label{tab:main-benchmarks}
\resizebox{\textwidth}{!}{\begin{tabular}{lllrrrr}
\toprule
패널 & 승식 & 기준모형 & 경주 수 & 개선폭 중앙값 & 95\% CI & 우위 판정 \\
\midrule
A & win & harville & 3,321 & -0.0679 & [-0.0691, -0.0666] & no \\
A & win & permutation & 3,321 & 0.2651 & [0.2610, 0.2686] & yes \\
A & win & other\_race & 3,320 & 0.3892 & [0.3843, 0.3939] & yes \\
A & win & uniform & 3,321 & 0.2567 & [0.2535, 0.2602] & yes \\
A & exacta & harville & 3,321 & 0.0573 & [0.0563, 0.0585] & yes \\
A & exacta & permutation & 3,321 & 0.4209 & [0.4162, 0.4245] & yes \\
A & exacta & other\_race & 3,320 & 0.5226 & [0.5174, 0.5289] & yes \\
A & exacta & uniform & 3,321 & 0.3758 & [0.3728, 0.3793] & yes \\
A & quinella & harville & 3,321 & 0.0464 & [0.0454, 0.0478] & yes \\
A & quinella & permutation & 3,321 & 0.3684 & [0.3640, 0.3719] & yes \\
A & quinella & other\_race & 3,320 & 0.4930 & [0.4882, 0.4982] & yes \\
A & quinella & uniform & 3,321 & 0.3444 & [0.3408, 0.3479] & yes \\
A & trio & harville & 3,321 & 0.1076 & [0.1061, 0.1094] & yes \\
A & trio & permutation & 3,321 & 0.4461 & [0.4434, 0.4491] & yes \\
A & trio & other\_race & 3,320 & 0.5532 & [0.5485, 0.5574] & yes \\
A & trio & uniform & 3,321 & 0.3988 & [0.3962, 0.4013] & yes \\
B & win & harville & 19,284 & -0.0862 & [-0.0868, -0.0857] & no \\
B & win & permutation & 19,284 & 0.3246 & [0.3230, 0.3265] & yes \\
B & win & other\_race & 19,283 & 0.4782 & [0.4756, 0.4808] & yes \\
B & win & uniform & 19,284 & 0.3213 & [0.3201, 0.3226] & yes \\
B & exacta & harville & 19,284 & 0.0494 & [0.0489, 0.0499] & yes \\
B & exacta & permutation & 19,284 & 0.5106 & [0.5089, 0.5127] & yes \\
B & exacta & other\_race & 19,283 & 0.6263 & [0.6242, 0.6285] & yes \\
B & exacta & uniform & 19,284 & 0.4661 & [0.4646, 0.4676] & yes \\
B & quinella & harville & 19,284 & 0.0353 & [0.0348, 0.0359] & yes \\
B & quinella & permutation & 19,284 & 0.4526 & [0.4508, 0.4544] & yes \\
B & quinella & other\_race & 19,283 & 0.5926 & [0.5906, 0.5945] & yes \\
B & quinella & uniform & 19,284 & 0.4254 & [0.4240, 0.4270] & yes \\
B & trio & harville & 19,284 & 0.1057 & [0.1049, 0.1064] & yes \\
B & trio & permutation & 19,284 & 0.5537 & [0.5519, 0.5550] & yes \\
B & trio & other\_race & 19,283 & 0.6589 & [0.6572, 0.6607] & yes \\
B & trio & uniform & 19,284 & 0.4943 & [0.4929, 0.4955] & yes \\
\bottomrule
\end{tabular}
}
\end{table}

표~\ref{tab:main-benchmarks}의 여러 기준모형 행과 사전 민감도 기준에 대해
별도의 다중비교 보정은 적용하지 않는다. 기준모형 개선 비교 32개와 세 민감도
기준에서의 두 패널 판정 24개는 모두 분석 전에 고정된 비교이며, 탐색 후 유의한
항목만 선택한 집합이 아니다. 여기서 목적은 가족단위 유의성 선언보다 각 사전 비교의
크기와 방향을 투명하게 보고하는 데 있으므로 별도 family-wise 보정은 적용하지 않는다.
따라서 각 95\% 신뢰구간을 가족단위 오류율이 통제된 동시 신뢰진술로 해석하지 않으며,
주결론은 사전에 정한 두 패널의 방향 일치와 개선폭의 경제적 크기를 함께 사용해
제한적으로 해석한다. 또한 타경주(``other\_race'') 기준은 동일 출전두수 안에서
결정론적으로 donor를 고르므로 얇은 출전두수 계층에서는 같은 donor 가격벡터가
여러 대상경주에 재사용될 수 있다. 따라서 표의 other\_race 95\% 신뢰구간은 donor
재사용을 별도 군집보정하지 않은 경주단위 bootstrap 개별구간이며, 이 기준은 주모형의
우월성을 확인하는 보조적 benchmark로 해석한다. 동결된
\texttt{outputs/main\_other\_race\_donor\_reuse.csv}에서 targets-per-distinct-donor는
Panel A의 출전두수별 계층에서 1.50--1.65, Panel B에서 1.42--1.68이고,
최대 재사용 횟수는 각각 2--6회와 3--7회다. 이 진단은 donor 의존성의 실제 크기를
가시화하지만 공동 주판정이나 사전 bootstrap 절차를 변경하지 않는다.

\begin{table}[htbp]
\centering
\caption{Panel A의 쌍승--복승 순서정보 비교}
\label{tab:main-order-a}
\resizebox{0.86\textwidth}{!}{\begin{tabular}{lrrrrrr}
\toprule
거리 & 경주 수 & 쌍승 개선 & 복승 개선 & 차이 중앙값 & 95\% CI 하한 & 95\% CI 상한 \\
\midrule
tv & 3,338 & 0.0573 & 0.0464 & 0.0105 & 0.0098 & 0.0111 \\
js & 3,338 & 0.0077 & 0.0063 & 0.0014 & 0.0013 & 0.0015 \\
\bottomrule
\end{tabular}
}
\end{table}

\begin{table}[htbp]
\centering
\caption{Panel B의 쌍승--복승 순서정보 부분식별}
\label{tab:main-order-b}
\resizebox{0.96\textwidth}{!}{\begin{tabular}{lrrrrrrl}
\toprule
거리 & 경주 수 & 차이 하한 & 하한 95\% CI & 차이 상한 & 상한 95\% CI & 보수적 양(+) 판정 \\
\midrule
tv & 19,284 & -0.0095 & [-0.0098, -0.0091] & 0.0314 & [0.0310, 0.0318] & no \\
\bottomrule
\end{tabular}
}
\end{table}

\begin{table}[htbp]
\centering
\caption{P3 공유 삼쌍승 가격집합을 강제한 출전두수별 1개 대표경주 sharpness 진단.
해 상태가 ``certified''인 행은 적어도 한 방향이 시간제한 때문에 sharp하지 않음을 뜻하며,
현재 16두 대표경주는 하한만 certified outer bound이고 상한은 sharp하다.}
\label{tab:main-order-joint}
\resizebox{\textwidth}{!}{\begin{tabular}{rrrrl}
\toprule
출전두수 & 기존 보수구간 & 공유가격 MILP 인증구간 & 폭 비율 & 해 상태 \\
\midrule
5 & [-0.0318, 0.0479] & [0.0025, 0.0133] & 0.135 & sharp \\
6 & [-0.0043, 0.0502] & [0.0176, 0.0265] & 0.164 & sharp \\
7 & [-0.0066, 0.0253] & [0.0038, 0.0148] & 0.346 & sharp \\
8 & [-0.0122, 0.0079] & [-0.0053, 0.0009] & 0.307 & sharp \\
9 & [-0.0487, -0.0096] & [-0.0347, -0.0234] & 0.291 & sharp \\
10 & [-0.0164, 0.0291] & [-0.0012, 0.0135] & 0.324 & sharp \\
11 & [-0.0192, 0.0138] & [-0.0082, 0.0019] & 0.304 & sharp \\
12 & [0.0114, 0.0316] & [0.0181, 0.0243] & 0.308 & sharp \\
13 & [0.0029, 0.0770] & [0.0262, 0.0482] & 0.297 & sharp \\
14 & [-0.0325, 0.0858] & [0.0079, 0.0359] & 0.236 & sharp \\
15 & [-0.1022, 0.1123] & [-0.0277, 0.0273] & 0.256 & sharp \\
16 & [-0.0556, 0.0550] & [-0.0203, 0.0107] & 0.280 & certified \\
\bottomrule
\end{tabular}
}
\end{table}

Panel A에서는 Harville 대비 개선폭이 쌍승과 복승 사이에서 유의하게 다르다.
하지만 Panel B의 보수적 차이 구간은 0을 사이에 둔다. 표~\ref{tab:main-order-joint}은
이 보수성이 쌍승과 복승에 서로 독립적인 삼쌍승 가격집합을 허용한 데서 얼마나
오는지 확인하기 위해, 출전두수 5--16두의 각 계층에서 결정론적으로 고른 대표경주
정확히 한 개씩에 대해 동일한 삼쌍승 가격변수가 두 하위 승식 주변분포를 동시에
생성하도록 강제한 joint MILP 결과를 보인다. 따라서 이 표는 경계의 느슨함을
진단하는 계산실험이지, 각 출전두수 계층에서 양·음 부호가 나타나는 경주의 비율이나
계층 평균효과를 추정한 표가 아니다. 인증구간의 폭은 기존 보수구간의 약
14--35\%로 크게 줄어든다. 5--7두와 12--14두의 선택된 대표경주 각각에서는
양(+)으로, 9두의 선택된 대표경주에서는 음(-)으로 부호가 식별되지만,
8·10·11·15·16두의 선택된 대표경주에서는 여전히 0을 포함한다. 16두의
최소화 문제는 제한시간 내 최적성을 완전히 인증하지 못해 그 하한은 certified
outer bound로 보고한다. 따라서 공유가격 제약을 정확히 반영하면 P3 경계가 상당히
날카로워질 수 있음을 확인하지만, 이 12개 진단 경주의 부호 패턴을 전체 시장이나
각 출전두수 계층의 부호분포로 일반화하지 않는다.

따라서 순서를 구분하는 쌍승에서 추가적인 가격정보가 더 강하다는 패턴은 clean
표본의 진단으로는 남지만, 전체 시장에 일반화되는 확정적 P3 결과로 보지 않는다.
이 결과는 처음부터 P3를 복합복권 비축약의 구조적 검정과 구분한 이유와도 일치한다.

종합하면, 삼쌍승식의 완전한 순서 가격벡터는 하위 복합승식의 가격을 단승 기반
Harville과 비정보적 기준보다 일관되게 더 잘 재구성한다. 동시에 Panel B는
배당상한과 반올림을 보수적으로 반영한 뒤에도 잔여 가격차가 경제적으로 0이라고
주장할 수 없음을 보여준다. 이후 행동모형 분석은 이 잔여 차이가 위험선호와
확률가중, 그리고 복합사건의 축약 방식과 어떤 관계를 갖는지를 별도의 식별
문제로 다룬다.
